\section{Подводные камни измерений и два отзыва}
\label{sec:pitfalls}

У оценок агентной памяти есть характерный отказ: числа реальны, скрипт
верен, а подпись на столбце — нет. Каждый пункт ниже стоил этому проекту
целой серии, прежде чем стать правилом. Мы их приводим, потому что правила
переносятся даже там, где наши результаты могут и не перенестись.

\subsection{Пять подводных камней}

\begin{enumerate}
\item \textbf{Самоотчёт модели — не вердикт.} Решает только собственный
      верификатор задачи. Каждая альтернатива, которую мы пробовали,
      производила строки, похожие на данные, но данными не являвшиеся.

\item \textbf{Экспериментальная рука может незаметно стать контролем.}
      Если плагин не загрузился, рука продолжает работать как базовая линия
      под другим именем, и числа её остаются правдоподобными. Поэтому гейт
      сверяет реализованный состав контекста с промптом, уходящим
      провайдеру, а не с собственным логом плагина\code.

\item \textbf{Детектор может поймать конфигурацию вместо событий.}
      Детектор тайм-аутов в этом проекте совпал со статической строкой
      политики повторов, которая печатается в каждом логе, и сообщил 54
      тайм-аута там, где их не было\live. Дефект проявился только через
      невозможное число ниже по течению: долю решённых задач выше
      \SI{100}{\percent}. Его модульный тест закодировал ту же ошибку, и
      набор тестов был зелёным, потому что код и тест ошибались одинаково.
      Принятое правило: прежде чем пускать детектор к корпусу, его
      проверяют на запуске с заранее известной истиной.

\item \textbf{Насыщение провайдера выглядит как отказ метода.} При
      конкурентности выше измеренного потолка провайдера \SI{68}{\percent}
      запросов получили ограничение скорости\live. Изголодавший запуск всё
      равно получает от верификатора честный ноль, и этот ноль попадает в
      таблицу как «метод провалился». Контаминация по построению
      однонаправленна: контроль уже завершился в чистых условиях, поэтому
      ходы теряли только экспериментальные руки. Правило карантина
      зафиксировали до того, как увидели подсчёты: чисто — ноль ошибочных
      ходов, номинально — не больше двух и меньше \SI{5}{\percent}, всё
      остальное исключается.

\item \textbf{Бинарные конечные точки без мощности.} Перед этой серией мы
      посчитали, что для обнаружения различия в доле решённых задач нужно
      около 72 задач при измеренной частоте несовпадений, и сказали это
      заранее, вместо того чтобы выдать нулевой результат за
      эквивалентность. Текущая серия это требование покрывает — именно
      поэтому \S\ref{sec:capability} вообще делает утверждение о
      решаемости.
\end{enumerate}

\subsection{Два отзыва}

Оба приводятся, потому что отзыв — тоже результат об измерении.

\paragraph{«Цензурирование по лимиту времени объясняет разрыв в
решаемости»: отозвано, цензурирования не существовало.} Это было следствие
дефекта детектора из пункта~3. После исправления детектора истинное число
цензурированных по времени запусков в корпусе оказалось нулевым. Последующее
наблюдение, что поднятие лимита вернуло четыре из семи провалившихся задач,
проверили против измеренной базовой частоты восстановления вердикта между
неизменёнными повторными запусками (\num{5}/\num{11}); вероятность увидеть
четыре или больше восстановлений только за счёт этой базовой частоты —
\num{0.40}\live. Ничего примечательного, отозвано.

\paragraph{«Только состояние менее стабильно»: отозвано на третьем
испытании.} После двух испытаний рука «только состояние» показывала более
высокую дисперсию между испытаниями на 10 из 12 задач
($p={}$\num{0.039}). Третье испытание сдвинуло это к 8 из 12
($p={}$\num{0.388}). Утверждение умерло ровно так, как и должен умереть
двухиспытательный $p$ чуть ниже \num{0.05}. Остановись серия на двух
испытаниях, в статью попал бы ложный эффект.

\paragraph{Чего отзывы не задели.} Каждое измерение стоимости, токенов и
кэша в этой статье не зависит от вердиктов, а результаты по решаемости из
\S\ref{sec:capability} получены дизайном, который появился позже всего
перечисленного выше и покрывает требование к мощности, его мотивировавшее.

\subsection{Промежуточный случай контаминации, оставленный как
предупреждение}

Ранний анализ схемы в этом проекте дал $p={}$\num{7.9e-08} в пользу
сильного эффекта. При повторном прогоне того же сравнения тем же скриптом
на карантинно-чистом корпусе получилось $p={}$\num{9.2e-02}\live.
Контаминированные ячейки были сосредоточены именно в сравнивавшихся руках:
арифметика была верна, а сбивала с толку подпись столбца.
